% Folha de prova das 20 figuras, na página A4 real das apostilas.
%
% Usa a própria classe `doeapostila`, e não um `article` qualquer, para que
% a largura de texto, a fonte e o corpo sejam exatamente os do caderno final.
% Cada figura entra na largura recomendada em CATALOGO.md: o que se lê aqui
% é o que o aluno vai ler na página impressa.
%
% Compilar (a classe está em ../estilo, fora deste diretório):
%   $env:PATH = "C:\Users\Narcísius\AppData\Local\Programs\MiKTeX\miktex\bin\x64;" + $env:PATH
%   pdflatex -interaction=batchmode -include-directory=..\estilo contato-visual.tex
\documentclass{doeapostila}
% ---------------------------------------------------------------------------
% macros.tex — notação matemática comum aos cinco cadernos.
% Fixada antes de qualquer caderno começar, para os cinco não divergirem.
% Uso: % ---------------------------------------------------------------------------
% macros.tex — notação matemática comum aos cinco cadernos.
% Fixada antes de qualquer caderno começar, para os cinco não divergirem.
% Uso: % ---------------------------------------------------------------------------
% macros.tex — notação matemática comum aos cinco cadernos.
% Fixada antes de qualquer caderno começar, para os cinco não divergirem.
% Uso: \input{../estilo/macros} no preâmbulo de cada main.tex.
% ---------------------------------------------------------------------------

% Somas de quadrados
\newcommand{\SQT}{\mathrm{SQ}_{\mathrm{T}}}    % total
\newcommand{\SQR}{\mathrm{SQ}_{\mathrm{R}}}    % da regressão
\newcommand{\SQr}{\mathrm{SQ}_{\mathrm{r}}}    % residual
\newcommand{\SQfaj}{\mathrm{SQ}_{\mathrm{faj}}}% falta de ajuste
\newcommand{\SQep}{\mathrm{SQ}_{\mathrm{ep}}}  % erro puro

% Médias quadráticas: \MQ{R}, \MQ{r}, \MQ{faj}, \MQ{ep}
\newcommand{\MQ}[1]{\mathrm{MQ}_{\mathrm{#1}}}

% Coeficientes de determinação
\newcommand{\Rdois}{R^{2}}
\newcommand{\Rajus}{R^{2}_{\mathrm{aj}}}

% Efeitos: \efeito{1}, \efeito{12}
\newcommand{\efeito}[1]{E_{#1}}

% Álgebra matricial do ajuste por mínimos quadrados
\newcommand{\vetbeta}{\boldsymbol{\beta}}
\newcommand{\matX}{\mathbf{X}}
\newcommand{\vety}{\mathbf{y}}
 no preâmbulo de cada main.tex.
% ---------------------------------------------------------------------------

% Somas de quadrados
\newcommand{\SQT}{\mathrm{SQ}_{\mathrm{T}}}    % total
\newcommand{\SQR}{\mathrm{SQ}_{\mathrm{R}}}    % da regressão
\newcommand{\SQr}{\mathrm{SQ}_{\mathrm{r}}}    % residual
\newcommand{\SQfaj}{\mathrm{SQ}_{\mathrm{faj}}}% falta de ajuste
\newcommand{\SQep}{\mathrm{SQ}_{\mathrm{ep}}}  % erro puro

% Médias quadráticas: \MQ{R}, \MQ{r}, \MQ{faj}, \MQ{ep}
\newcommand{\MQ}[1]{\mathrm{MQ}_{\mathrm{#1}}}

% Coeficientes de determinação
\newcommand{\Rdois}{R^{2}}
\newcommand{\Rajus}{R^{2}_{\mathrm{aj}}}

% Efeitos: \efeito{1}, \efeito{12}
\newcommand{\efeito}[1]{E_{#1}}

% Álgebra matricial do ajuste por mínimos quadrados
\newcommand{\vetbeta}{\boldsymbol{\beta}}
\newcommand{\matX}{\mathbf{X}}
\newcommand{\vety}{\mathbf{y}}
 no preâmbulo de cada main.tex.
% ---------------------------------------------------------------------------

% Somas de quadrados
\newcommand{\SQT}{\mathrm{SQ}_{\mathrm{T}}}    % total
\newcommand{\SQR}{\mathrm{SQ}_{\mathrm{R}}}    % da regressão
\newcommand{\SQr}{\mathrm{SQ}_{\mathrm{r}}}    % residual
\newcommand{\SQfaj}{\mathrm{SQ}_{\mathrm{faj}}}% falta de ajuste
\newcommand{\SQep}{\mathrm{SQ}_{\mathrm{ep}}}  % erro puro

% Médias quadráticas: \MQ{R}, \MQ{r}, \MQ{faj}, \MQ{ep}
\newcommand{\MQ}[1]{\mathrm{MQ}_{\mathrm{#1}}}

% Coeficientes de determinação
\newcommand{\Rdois}{R^{2}}
\newcommand{\Rajus}{R^{2}_{\mathrm{aj}}}

% Efeitos: \efeito{1}, \efeito{12}
\newcommand{\efeito}[1]{E_{#1}}

% Álgebra matricial do ajuste por mínimos quadrados
\newcommand{\vetbeta}{\boldsymbol{\beta}}
\newcommand{\matX}{\mathbf{X}}
\newcommand{\vety}{\mathbf{y}}

\definecaderno{--}{Figuras}

\newcommand{\prova}[4]{%
  % #1 caderno · #2 arquivo · #3 fração de \linewidth · #4 o que ensina
  \begin{figure}[H]
    \centering
    \includegraphics[width=#3\linewidth]{out/#2}
    \caption{\texttt{#2} \quad·\quad caderno #1 \quad·\quad
             \texttt{width=#3\textbackslash linewidth}\\ #4}
  \end{figure}}

\begin{document}

\begin{center}
  {\large\bfseries Folha de prova das figuras}\par\smallskip
  {\small Planejamento de Experimentos · 20 figuras · página A4 real}\par
\end{center}

\noindent
Esta folha existe para uma única pergunta: cada figura continua legível no
tamanho em que será impressa? Todas aparecem abaixo na largura recomendada em
\texttt{CATALOGO.md}, sobre a mesma mancha de texto dos cinco cadernos.
Rótulo de eixo, legenda e anotação devem ser lidos sem esforço a 100\,\% de
zoom. Preto e branco, com uma única cor de apoio, o azul \texttt{\#1B4E79}.

\section*{Caderno 1 — Por que planejar experimentos}

\prova{1}{c1-um-fator-por-vez}{0.80}{%
  A busca um fator de cada vez estaciona longe do ótimo quando há interação.}
\prova{1}{c1-dispersao-20-medidas}{1.00}{%
  Vinte medidas contra a especificação de projeto: a tabela não conclui, o
  gráfico sim.}
\prova{1}{c1-blocagem-intervalo}{1.00}{%
  Os mesmos dez valores levam a conclusões opostas conforme se reconheça ou
  não o fabricante como bloco.}
\prova{1}{c1-variancia-da-media}{0.70}{%
  $\mathrm{Var}(\bar{x}) = \sigma^{2}/N$: o retorno decrescente da réplica.}

\section*{Caderno 2 — Planejamento fatorial $2^{k}$}

\prova{2}{c2-cubo-2a3}{0.60}{%
  O $2^{3}$ com as respostas nos vértices e o efeito principal lido como
  diferença entre faces opostas.}
\prova{2}{c2-interacao}{1.00}{%
  Retas que se cruzam indicam interação; retas paralelas, ausência dela.}
\prova{2}{c2-pareto-efeitos}{1.00}{%
  Pareto dos efeitos padronizados e o que a linha crítica separa.}
\prova{2}{c2-normal-efeitos}{0.75}{%
  Sem réplica, os efeitos que fogem da reta normal são os reais.}

\section*{Caderno 3 — Regressão e análise de variância}

\prova{3}{c3-nuvem-reta-residuos}{1.00}{%
  Nuvem, reta ajustada e os resíduos como segmentos verticais.}
\prova{3}{c3-decomposicao-sq}{1.00}{%
  $\SQT = \SQR + \SQr$, ponto a ponto e em proporção.}
\prova{3}{c3-banda-confianca}{0.60}{%
  A banda de previsão em ampulheta: mais estreita em $\bar{x}$, aberta nas
  bordas.}
\prova{3}{c3-diagnosticos-residuo}{1.00}{%
  Os quatro diagnósticos de resíduo sobre um modelo com falta de ajuste.}

\section*{Caderno 4 — Curvatura e superfície de resposta}

\prova{4}{c4-grade-2a2-ponto-central}{0.95}{%
  A grade $2^{2}$ com ponto central, em unidades naturais e codificadas.}
\prova{4}{c4-linear-vs-quadratico}{0.80}{%
  Os mesmos dezoito pontos sob o ajuste linear e sob o quadrático.}
\prova{4}{c4-superficie-plana-vs-curva}{0.90}{%
  Um modelo de primeira ordem não tem máximo interior; um com termos
  quadráticos tem.}
\prova{4}{c4-composicao-central}{0.60}{%
  Geometria da composição central rotacionável de dois fatores.}
\prova{4}{c4-contorno-otimo}{0.80}{%
  O ponto ótimo da superfície não coincide com o melhor ensaio realizado.}

\section*{Caderno 5 — Economia de ensaios}

\prova{5}{c5-projecao-meia-fracao}{1.00}{%
  Uma fração com fator inerte volta a ser fatorial completo, com réplicas.}
\prova{5}{c5-ensaios-por-arranjo}{0.75}{%
  Número de ensaios por arranjo, em escala logarítmica.}
\prova{5}{c5-corpos-de-prova-vs-informacao}{1.00}{%
  Com orçamento fixo de material, mais fatores custam graus de liberdade.}

\end{document}
